\nuevaReceta{Chocolate a la Taza}{1 taza por persona}{10-15 min}{receta:chocolate_taza}

    \begin{minipage}[t]{0.35\textwidth}
        \section*{Ingredientes}
        \begin{itemize}[leftmargin=*]
            \item 1 pastilla de Chocolate Valor a la Taza
            \item 1 taza de leche por persona
            \item 3-5 onzas de chocolate por cada taza de leche
        \end{itemize}
    \end{minipage}
    \hfill
    \begin{minipage}[t]{0.6\textwidth}
        \section*{Preparación}
        \begin{enumerate}
            \item \textbf{Calcular el chocolate:} Separar entre 3 y 5 onzas de la pastilla por cada taza de leche. Con 3 quedará más ligero; con 5, bastante espeso. El punto habitual de esta receta está en 4 o 4,5 onzas por taza.
            \item \textbf{Calentar la leche:} Poner la leche en un cazo a fuego medio y calentarla sin dejar que llegue a hervir.
            \item \textbf{Añadir el chocolate:} Incorporar las onzas y remover continuamente, procurando que no se pegue al fondo.
            \item \textbf{Espesar:} Mantener a fuego medio-bajo hasta que el chocolate se haya fundido por completo y alcance el espesor deseado.
            \item \textbf{Servir:} Repartir inmediatamente en las tazas. Al enfriarse espesará un poco más.
        \end{enumerate}
    \end{minipage}

    \vspace{1em}
    \begin{tcolorbox}[colback=orange!5, colframe=orange!50, title=El punto del chocolate, size=small]
        La cantidad de onzas marca el espesor final: 3 para un chocolate más ligero, 4-4,5 para el punto habitual de casa y 5 para uno muy espeso. Si queda demasiado denso, añade un poco de leche caliente; si queda ligero, incorpora otra media onza y sigue removiendo.
    \end{tcolorbox}
\end{receta}
